\section{Нулевая базовая кривизна: структура и конечная остановка}
\label{sec:zero}
Предельный случай $H_0=0$ качественно отличается от малого
положительного $H_0$. Утверждения ниже относятся к условному зазору
$\delta=f-f_Q^*$; условие кривизны выполняется на всём $Q$.

\begin{proposition}[Экспонированная оптимальная грань и общий градиент]
\label{prop:zeroface}
Пусть $H_0=0$. Для любого $x^*\in X^*:=\arg\min_Q f$ положим
$g^*=\nabla f(x^*)$. Тогда
\begin{equation}\label{eq:zeroface}
 X^*=\{x\in Q:\inner{g^*}{x-x^*}=0\}.
\end{equation}
Во всех оптимальных точках градиент равен одному и тому же $g^*$,
а для $x\in Q$
\begin{equation}\label{eq:zerograd}
 \norm{\nabla f(x)-g^*}\le\tfrac12H_1D\delta(x).
\end{equation}
В частности, стационарный в объемлющем пространстве оптимум $g^*=0$
вынуждает $f$ быть постоянной на $Q$. Непостоянная функция этого
класса с нулевой базовой кривизной также не может иметь оптимум в
относительной внутренности.
\end{proposition}
\begin{proof}
Оптимальность и выпуклость дают $w:=\inner{g^*}{x-x^*}\ge0$ и
$\delta(x)\ge w$. Для $h(t)=\delta(x^*+t(x-x^*))$ условие
кривизны даёт $h''\le H_1\norm{x-x^*}^2h$, причём $h(0)=0$ и
$h'(0)=w$. То же сравнение с положительным ядром, что и в
теореме~\ref{lem:bridge}, даёт
\[
 0\le h(t)\le w\frac{\sinh(rt)}r,\qquad
 r=\sqrt{H_1}\norm{x-x^*},
\]
где при $r=0$ отношение равно $t$. Поэтому $w=0$ влечёт
$h(1)=0$, что доказывает~\eqref{eq:zeroface}.
Отрезок между двумя оптимумами лежит в $X^*$; гессиан на нём
нулевой по предполагаемому ограничению нормы, поэтому интегрирование
показывает совпадение градиентов. На отрезке от $x^*$ до $x$
выпуклость даёт $h(t)\le t\delta(x)$. Интегрируя норму гессиана,
получаем
\[
 \norm{\nabla f(x)-g^*}
 \le H_1\norm{x-x^*}\int_0^1h(t)\,dt
 \le\tfrac12H_1D\delta(x).
\]
В оптимуме из относительной внутренности линейный функционал $g^*$
постоянен на аффинной оболочке $Q$, поэтому~\eqref{eq:zeroface}
даёт $X^*=Q$.
\end{proof}

Пусть теперь $Q=\operatorname{conv}(\mathcal A)$ — многогранник
с конечным множеством атомов, а $f$ непостоянна на $Q$.
Обозначим $\mathcal A_*:=\mathcal A\cap X^*$ и
\begin{equation}\label{eq:atommargin}
 m:=\min_{v\in\mathcal A\setminus\mathcal A_*}
              \inner{g^*}{v-x^*}>0.
\end{equation}
Для любого активного представления $x=\sum\lambda_vv$ масса вне
оптимальной грани удовлетворяет
\begin{equation}\label{eq:badmass}
 \theta:=\sum_{v\notin\mathcal A_*}\lambda_v
 \le\frac{\inner{g^*}{x-x^*}}m\le\frac{\delta(x)}m.
\end{equation}
Замена неоптимальных атомов любой точкой $X^*$ показывает, что
$\operatorname{dist}(x,X^*)\le D\theta$. Следовательно, сама
модель влечёт глобальную оценку слабой остроты
\[
 \delta(x)\ge(m/D)\operatorname{dist}(x,X^*).
\]
Независимая гипотеза квадратичного роста здесь не нужна.
Конечность множества атомов существенна: для аффинной функции
$f(x)=1-x_1$ на единичном диске
\[
\delta(\cos t,\sin t)/\operatorname{dist}((\cos t,\sin t),X^*)
=|\sin(t/2)|\to0.
\]
Таким образом, нулевая базовая кривизна на
произвольном компактном выпуклом множестве не влечёт слабой остроты.

\begin{theorem}[Конечная точная остановка AFW с обратным поиском]
\label{thm:zerofinite}
Пусть $H_0=0$, множество атомов конечно и $U_0>0$. Положим
\[
 b=H_1D^2,\quad
 c_b=\begin{cases}\sqrt b\coth\sqrt b,&b>0,\\1,&b=0,\end{cases}
 \qquad \overline L=\max\{M,2H_1U_0\},
\]
\begin{equation}\label{eq:zerothreshold}
 \tau=\min\left\{\frac{U_0}2,
       \frac{m}{4(1+b+c_b)}\right\},\qquad
 J_0=\ceil{\log_2(U_0/\tau)}.
\end{equation}
Определим два потенциала инициализации:
\[
 P=\ceil{\log_2\max\{1,MD^2/U_0\}},\qquad
 Q_0=\ceil{\log_2\max\{1,8\overline L D^2\tau/m^2\}}.
\]
Стандартный алгоритм AFW с обратным поиском, продолжаемый с нулевым
допуском остановки, достигает точного оптимума и $U_N=0$ за
\begin{equation}\label{eq:zerofinite}
 N\le2\big[P+\ceil{(6+16b)J_0}+Q_0\big]+|S_0|-1
\end{equation}
обновлений. После первой исходной точки с $U\le\tau$ каждое
неоптимальное обновление выбирает направление удаления; до момента,
когда остаются лишь максимальные удаления неоптимальных атомов,
происходит не более $Q_0$ дополнительных хороших обновлений.
Числа градиентов и полных LMO равны $N+1$, а
\[
 N_f\le1+2N+\ceil{\log_2\max\{1,2H_1U_0/M\}}.
\]
Ни $m,H_1,f_Q^*$, ни порог $\tau$ не являются входами алгоритма.
Не требуется ни гипотеза~\eqref{eq:geometry}, ни положительная
константа сильной выпуклости. Утверждения предполагают точную арифметику.
\end{theorem}
\begin{proof}
Огибающая принятой кривизны даёт $L_k\le\overline L$ для всех
обновлений. Рассмотрим неоптимальную исходную точку с $U\le\tau$;
тогда $0<\delta\le\tau$. По~\eqref{eq:zerograd} различие между
текущей и оптимальной разностями градиентных оценок любых двух атомов
не превосходит $b\delta/2$. Поэтому текущая оценка любого
неоптимального атома строго больше оценки любого оптимального.
LMO возвращает оптимальный атом. У неоптимальной исходной точки
есть активный неоптимальный атом, и максимизатор удаления $v$
обязательно неоптимален. Его зазор удовлетворяет
\[
 g_{\mathrm A}=\inner{\nabla f(x)}{v-x}
 \ge m-\delta-b\delta/2\ge3m/4.
\]
В то же время теорема~\ref{lem:bridge} даёт
$G\le c_b\delta\le m/4$. Поэтому алгоритм выбирает удаление.
Вес этого атома не превосходит $\theta\le\delta/m<1/2$, откуда
\[
 \alpha_{\max}\le2\delta/m.
\]
Положим $R=m^2/(8D^2\tau)$. Поскольку
\[
 2H_1\tau/R=16b\tau^2/m^2
 \le b/(1+b+c_b)^2\le1,
\]
принятая оценка выше $R$ обязательно получена делением пополам на
первой пробе; приёмка после отказа даёт менее
$2H_1\delta\le R$. При входе унаследованная оценка не превосходит
$\overline L$, поэтому таких делений не более $Q_0$.
Принятая оценка не выше $R$ остаётся не выше $R$ по локальной
огибающей и удовлетворяет
\[
 L_k\norm d^2\alpha_{\max}
 \le 2RD^2\delta/m\le m/4<g_{\mathrm A}.
\]
Каждая такая проба обрезается на допустимой верхней границе и удаляет
неоптимальный атом. Целевая функция и $U$ убывают, поэтому локальные
условия сохраняются. После конечного числа удалений все активные
атомы оптимальны. Тогда $G=0$, и оплаченная финальная проверка даёт
$U=0$.

Для подсчёта предшествующих обновлений применим фазовое доказательство
теоремы~\ref{thm:afw} до первой исходной точки с $U\le\tau$.
Оно ограничивает суммарное число хороших обновлений до этой точки
величиной $P+\ceil{(6+16b)J_0}$. Локальный аргумент допускает
не более $Q_0$ дополнительных хороших обновлений. Глобальный подсчёт
удалений~\eqref{eq:dropcount} доказывает~\eqref{eq:zerofinite}.
Тождество числа вызовов функции и огибающая дают заявленную
оракульную оценку, включая отказы.
\end{proof}
Здесь $J_0$ не зависит от $M$. Поэтому для фиксированных функции
и старта большое $M$ даёт лишь $O(\log M)$ через $P,Q_0$, без
дополнительного множителя $H_1D^2$ при этой цене инициализации.

Конечная идентификация активного множества при строгой
дополняемости — известная теория~\cite[теорема 3.1 и приложение]{bomze},
а не новое общее явление. Garber~\cite[теоремы 4--5]{garber} также
даёт результаты конечной сходимости и идентификации грани для
вариантов FW с линейным поиском. Здесь сама модель с нулевой базовой
кривизной вынуждает экспонированную грань совпадать со всем множеством
решений; идентификация этой грани точно решает задачу. Приведённые
порог и число проб специализируют механизм к унаследованной оценке
кривизны с делением и умножением на два, не требуя сходимости к
конкретному оптимуму. Приоритет перед всеми теоремами идентификации
не заявляется.

\paragraph{Почему предельный случай существенен.}
Вывод не переносится на произвольное $H_0>0$. На $[-1,1]$
возьмём $f(x)=h x^2/2$, $h>0$, $0<x_0<1$ и $M=3h$.
Обратный поиск на каждом обновлении принимает $L_k=3h/2$:
начальная первая проба принимается, а каждая последующая проба
$3h/4$ отклоняется для этого ненулевого квадратичного шага.
Выбранное направление — FW в обычном представлении через два
конца отрезка, и
\[
 x_{k+1}=x_k-f'(x_k)/(3h/2)=x_k/3.
\]
Поэтому ни одна конечная итерация не оптимальна, сколь бы мало ни
было $h=H_0>0$. Конечная идентификация оптимальной грани в общем
случае не решает непостоянную на этой грани целевую функцию.
